\documentclass[11pt]{article}

\usepackage[T1]{fontenc}
\usepackage{lmodern}
\usepackage[margin=0.9in]{geometry}
\usepackage{microtype}
\usepackage{amsmath,amssymb}
\usepackage{booktabs,tabularx,array}
\usepackage{xcolor}
\usepackage{hyperref}

\definecolor{paperblue}{HTML}{17365D}
\hypersetup{
  colorlinks=true,
  linkcolor=paperblue,
  citecolor=paperblue,
  urlcolor=paperblue,
  pdftitle={Terminus 1.5B: A Peri-Normalized Value-Residual Transformer for Dense-to-Ternary Training},
  pdfauthor={Compact Intelligence}
}
\setlength{\parindent}{0pt}
\setlength{\parskip}{0.55em}
\renewcommand{\arraystretch}{1.15}
\newcommand{\model}{Terminus 1.5B}
\newcommand{\QAT}{quantization-aware training}

\title{\vspace{-1.0cm}\textbf{Terminus 1.5B: A Peri-Normalized Value-Residual Transformer for Dense-to-Ternary Training}}
\author{Compact Intelligence}
\date{Technical design paper \quad 14 October 2026}

\begin{document}
\maketitle

\begin{abstract}
We present the selected architecture and training method for \model{}, a
1,500,456,496-parameter decoder-only transformer intended for a staged
dense-to-ternary language-model program. The model combines a 4:1 local/global
attention schedule, grouped-query attention, Peri-LN branch normalization,
learned Q/K normalization, a cross-layer learned value residual, and a tied
65,536-token embedding and readout. Dense training uses FP8 E4M3 projection
arithmetic with FP32 master projection weights and spectral Muon updates;
ternary-weight continuation is a separate \QAT{} phase. The program targets
one trillion valid training-token presentations across both stages and
evaluates Hindi, Bengali, Tamil, Telugu, Kannada, Marathi, and Gujarati as its
initial Indic-language support set. We give the exact model geometry,
quantizer and optimizer rules, parameter and storage accounting, evaluation
protocol, and limits of the existing evidence. Production training, final
ternary quality, distributed throughput, and on-device performance are not
reported as completed results.
\end{abstract}

\noindent\textbf{Keywords:} small language model; dense-to-ternary training;
FP8; grouped-query attention; Peri-LN; Indic languages; offline inference

\section{Introduction}

Offline language assistance imposes two distinct constraints. A model must
learn useful multilingual behavior, and its exported form must fit the memory,
latency, and energy limits of its target device. A small parameter count alone
does not establish either property. Ternary storage can reduce the byte cost
of selected projection matrices, but a tied vocabulary table, scale metadata,
KV caches, activations, and runtime workspaces remain. Likewise, fast FP8
training arithmetic does not establish the quality of a later ternary model.

\model{} therefore follows a two-stage method: dense-weight pretraining with
FP8 computation, followed by ternary-weight \QAT{}. This sequence permits
learning in the dense model before testing how much function survives
conversion and adaptation. The selected architecture adds Peri-LN-style
normalization of residual branches and a learned value path from the first
attention layer to later layers. Both mechanisms are motivated by residual
scale and attention-value behavior observed in earlier proxy work; their
full-scale effects must be evaluated on the trained model
\cite{peri,valueresidual}.

This paper is a technical design and measurement protocol. It identifies
implemented model and optimizer choices, the stated one-trillion-token
training target, and the evidence needed for release claims. It contains no
claim that the production model has already processed that budget or reached
a specified multilingual benchmark score.

\section{Model architecture}

\subsection{Geometry and parameter allocation}

The model is a causal decoder with 25 blocks. Its hidden width is 2,304 and
the SwiGLU intermediate width is 5,760. Each attention layer has 18 query
heads and six key/value heads, with head dimension 128 and a uniform 3:1
query-to-KV ratio. The tied input embedding and output classifier use a frozen
65,536-token vocabulary. The starting sequence length is 4,096 tokens.

\begin{table}[ht]
\centering
\caption{Selected model and dense-stage implementation. Values describe the
selected configuration, not a completed production checkpoint.}
\label{tab:model}
\small
\begin{tabularx}{\linewidth}{@{}>{\bfseries}p{0.39\linewidth}X@{}}
\toprule
Property & Selection \\
\midrule
Trainable parameters & 1,500,456,496 \\
Decoder layers; hidden width & 25; 2,304 \\
SwiGLU intermediate width & 5,760 \\
Query / KV heads; head dimension & 18 / 6; 128 \\
Vocabulary and readout & 65,536 tokens; input and output weights tied \\
Starting context; local window & 4,096 tokens; 1,024 tokens \\
Attention schedule & Four local layers followed by one global layer, repeated five times \\
Normalization & Pre-attention and pre-FFN RMSNorm; Peri-LN branch-output RMSNorm; learned Q/K RMSNorm \\
Value path & Learned mixture of layer-zero and current-layer values in layers 2--25 \\
Dense linear computation & Block-scaled FP8 E4M3 through DeepGEMM; full FP8 projection backward \\
Residual stream; attention kernel & BF16; FlashAttention-3 where supported \\
Projection optimizer & Spectral Muon on complete logical matrices \\
\bottomrule
\end{tabularx}
\end{table}

For this configuration, 1,349,222,400 parameters belong to the projection
matrices eligible for later ternary storage. The tied embedding/readout has
150,994,944 parameters; other learned parameters total 239,152. These three
counts sum exactly to 1,500,456,496. The tied table is not part of the
present ternary projection set.

\subsection{Local and global attention}

Twenty layers use causal sliding-window attention over 1,024 tokens; every
fifth layer uses full causal attention over the 4,096-token starting context.
This four-local/one-global pattern gives bounded local-layer KV history and
periodic mixing across the full context. The head geometry and GQA ratio are
identical in both layer types. Grouped-query attention reduces KV-cache
storage relative to one KV head per query head \cite{gqa}. The chosen 4,096
context is an initial training specification; longer-context behavior
requires a separate training and evaluation record.

Query and key vectors are normalized with learned gains before rotary
position encoding. Values are not Q/K-normalized. The selected path has no
post-optimizer Q/K rescaling rule; global gradient clipping is part of the
optimizer configuration and serves a different purpose.

\subsection{Peri-LN residual blocks}

Write $R_{\gamma}(z)=\gamma\odot z/\sqrt{\operatorname{mean}(z^2)+\epsilon}$
for RMS normalization, with $\epsilon=10^{-5}$ and a learned gain. For block
$\ell$, let $x_\ell$ denote its input, and define
\begin{align}
u_\ell &= R_{\gamma^{\mathrm{in,a}}_\ell}(x_\ell),\\
a_\ell &= x_\ell +
  R_{\gamma^{\mathrm{out,a}}_\ell}
  \!\left(W^O_\ell\,\operatorname{Attn}_\ell(u_\ell)\right),\\
h_\ell &= R_{\gamma^{\mathrm{in,f}}_\ell}(a_\ell),\\
x_{\ell+1} &= a_\ell +
  R_{\gamma^{\mathrm{out,f}}_\ell}
  \!\left(W^D_\ell
  [\operatorname{SiLU}(W^G_\ell h_\ell)\odot(W^U_\ell h_\ell)]\right).
\end{align}
The two output normalizers act immediately before residual addition and
have gains independent of the input normalizers. This is the Peri-LN
component. The attention-output and FFN-down projection weights also receive
the selected initialization multiplier 0.14078571632817444. It is an
initialization rule, not a runtime multiplier. The model uses no SubLN
branch \cite{rmsnorm,swiglu,peri}.

\subsection{Cross-layer value residual}

Let $v_\ell$ be the value projection of the token at layer $\ell$. The first
layer keeps $\widetilde v_0=v_0$; every later layer uses
\begin{equation}
\widetilde v_\ell=\lambda_{\ell,0}v_0+
\lambda_{\ell,1}v_\ell,\qquad \ell\geq1.
\end{equation}
The two scalars per later layer are learned, initialized to 0.5 each, and
are not constrained to sum to one. Mixing occurs before attention masking
and aggregation. Training preserves the gradient path to the first value
projection. Cached inference maintains a rolling buffer of the first-layer
values for the new positions. The mechanism reuses low-layer information
without adding an MTP drafter or changing the attention cadence
\cite{valueresidual}.

\section{Dense-to-ternary training method}

\subsection{Dense FP8 stage}

The selected dense mode keeps FP32 master projection weights. Linear
operands use block-scaled FP8 E4M3 weights and activations through DeepGEMM;
the selected backward path uses FP8 projection arithmetic. The residual
stream is BF16. Attention uses FlashAttention-3 where supported, packed
Q/K/V and gate/up projections reduce launch and memory overhead, and the
tied output classifier uses streaming cross-entropy rather than an FP8
readout head \cite{fa3,deepgemm}. The training objective is ordinary
next-token cross-entropy on tokenized text streams.

The precise distinction matters: ``dense FP8'' describes arithmetic applied
to dense-weight operands, not a model whose only stored parameters are FP8.
Master projection weights and optimizer state retain higher precision.

\subsection{Ternary continuation}

In the currently implemented QAT path, each 128-by-128 projection-weight
tile $B$ has a statistical scale
\begin{equation}
s_B=\max\!\left(\frac{1}{|B|}\sum_{i,j\in B}|W_{ij}|,10^{-6}\right),
\qquad
Q_B(W)=s_B\operatorname{clip}
\!\left(\operatorname{round}(W_B/s_B),-1,1\right).
\end{equation}
The forward path encodes these scaled ternary values for FP8 matrix
multiplication; FP8 activations remain in use. A straight-through rule
passes projection-weight gradients to FP32 master weights. The scale is
recomputed from weights, so it is not a separate learned scale parameter.
The tied embedding/readout and learned normalization gains are outside the
ternary projection set. One scale per tile covers 16,384 weights, a
representation choice that must be assessed by final loss and tasks rather
than by code sparsity alone \cite{bitnet}.

The QAT transition changes the effective projection weights and can
produce an immediate loss increase. A measured transition report separates
loss at the conversion point from recovery during QAT and from final
post-QAT quality. Its learning rates, stable/decay duration, data mixture,
and optimizer-state policy are phase-specific. No dense-stage rate is
silently reused as a QAT prescription.

\subsection{Spectral Muon and auxiliary AdamW}

Packed Q/K/V and gate/up physical parameters are split into their logical
matrices for the Muon update. For a logical matrix $W$ with dimensions
$(d_{\mathrm{out}},d_{\mathrm{in}})$, let $U_t$ be the pinned five-step
Newton--Schulz direction calculated from its momentum-adjusted gradient.
The selected update is
\begin{equation}
W_{t+1}=(1-\eta_t\lambda)W_t-
\eta_t\sqrt{d_{\mathrm{out}}/d_{\mathrm{in}}}\,U_t.
\end{equation}
The aspect-ratio factor applies to each logical matrix, including each
slice of a packed projection. Weight decay uses the base $\eta_t$. This
rule has no additional width-growing Muon factor \cite{muon}.

\begin{table}[ht]
\centering
\caption{Selected dense-stage starting settings. Production batch and token
horizons are bound to the run manifest.}
\label{tab:optimizer}
\small
\begin{tabularx}{\linewidth}{@{}>{\bfseries}p{0.65\linewidth}X@{}}
\toprule
Setting & Value \\
\midrule
Muon base learning rate & 0.017 \\
Tied embedding/readout AdamW learning rate & 0.003 \\
Ordinary, Peri-LN, final-norm, value-mix AdamW learning rate & 0.0003 \\
Learned Q/K norm AdamW learning rate & 0.0001 \\
Hidden-matrix weight decay; tied-table decay & 0.02; 0 \\
Global gradient-norm clip & 1.0 \\
Schedule form; final learning-rate fraction & Warmup--stable--decay with cosine decay; 0.1 \\
\bottomrule
\end{tabularx}
\end{table}

These values are starting coordinates selected from smaller-width dense
research. A 200M-parameter campaign selected the matrix and small-group
rates; a matched 700M-width bridge favored the 0.003 embedding rate over
0.00424 and 0.006 at 500M tokens of English training. Those observations do
not establish a universal width law or a one-trillion-token optimizer
trajectory.

\section{Data scope and language coverage}

The program target is \textbf{1,000,000,000,000 valid training-token
presentations across the dense and QAT stages}. Padding and held-out
evaluation tokens do not count toward this total. Presented tokens and
distinct source tokens are separate quantities: any reuse is recorded as
exposure rather than described as new unique data. The production manifests
bind tokenizer identity, shards, data order, source proportions, and stage
boundaries. No fixed 40\% Indic allocation or inherited source-token table
is asserted for this revised program.

The initial Indic-language support and reporting set comprises \textbf{Hindi,
Bengali, Tamil, Telugu, Kannada, Marathi, and Gujarati}. English and
relevant code and mathematics are included in broader training and
evaluation. The seven-language set identifies intended coverage; it is not
a substitute for native-language benchmark results.

\begin{table}[ht]
\centering
\caption{Initial Indic-language evaluation set.}
\label{tab:languages}
\small
\begin{tabular}{@{}lll@{}}
\toprule
Language & ISO code & Primary script \\
\midrule
Hindi & hi & Devanagari \\
Bengali & bn & Bengali \\
Tamil & ta & Tamil \\
Telugu & te & Telugu \\
Kannada & kn & Kannada \\
Marathi & mr & Devanagari \\
Gujarati & gu & Gujarati \\
\bottomrule
\end{tabular}
\end{table}

The frozen 65,536-piece tokenizer has explicit Indic and English/code/math
handling. Corpus preparation requires language identification, script and
Unicode checks, near-duplicate removal, source-overlap accounting,
license/provenance review, sensitive-data filtering, and evaluation-set
decontamination. Synthetic and translated material must be labeled so that
native-language results cannot be inferred from translated-English scores.
The final data record reports per-language and per-source exposure,
tokenizer fertility, and train/validation boundaries. A claimed supported
language is evaluated individually, including relevant code-mixed use.

\section{Distributed execution and compute accounting}

The starting distributed strategy is replicated data parallelism on 32
H100-class ranks, with tensor, pipeline, and context parallel degrees of
one. Each rank retains the model and optimizer state. Accumulation keeps a
reviewed global valid-token batch, and DDP reduces FP32 gradients before
global clipping and the optimizer update. A microbatch of four 4,096-token
sequences is the first memory candidate, not a selected global batch.
Changing world size or microbatch must preserve the intended global batch
and numerical update contract.

For $W$ ranks, per-rank microbatch $m$, context $S=4096$, and accumulation
$a$, the nominal update has $B=WmSa$ valid tokens when all positions are
valid. The run configuration must derive $a$ from the chosen manifest and
global batch. A 32-rank system requires validation of rank-local data
ordering, all-reduce semantics, kernel parity, checkpoint recovery, and
full-update timing on its actual NVL/PCIe topology.

At a planning envelope of 6,000 aggregate GPU-hours, one trillion valid
tokens require
\begin{equation}
\bar r=\frac{10^{12}}{6000\times3600}
=46{,}296\ \text{tokens s}^{-1}\text{ GPU}^{-1}
\end{equation}
\emph{including} all allocated time consumed by dense training, QAT,
distributed communication, evaluation, checkpoints, and recovery. Thirty-two
GPUs would provide 187.5 wall-clock allocation hours under that envelope.
An optimized single-H100 NVL dense measurement averaged 40,996 tokens/s
over three complete 20-update receipts, with 55.5 GB peak allocated memory
at its measured batch. If that dense rate held without any overhead for all
6,000 GPU-hours, it would yield roughly 886B tokens. This arithmetic does
not establish cluster performance or the QAT rate. The one-trillion-token
target is therefore a program target whose schedule must be accepted from
measured full-cluster, all-in throughput, not from device peak FLOPs.

\section{Evidence and evaluation protocol}

\subsection{Evidence already available}

The prior 146,423,680-parameter proxy is informative about conversion risk,
but differs from the selected 1.5B architecture and optimizer. It completed
10B dense tokens and a 2B ternary-QAT continuation without a recorded
nonfinite incident. Under one fixed held-out weighting, its mean per-token
negative log-likelihood moved from 2.4194 at the dense checkpoint to 2.6042
at the final QAT checkpoint. The QAT training mixture also changed because
fresh English material was limited, so the final difference cannot be
attributed to quantization alone. At the exact conversion point, before any
QAT update, held-out loss rose sharply, demonstrating that the weight-mode
change perturbed model function. The proxy result motivates a controlled
QAT transition; it is not a prediction of 1.5B final quality.

Architecture selection, dense optimizer tuning, the older 146M run, and
future production training have different model or optimizer contracts.
Their losses are not pooled into a single comparison. Reported numerical
results state the checkpoint, data mixture, scorer, token count, and
precision path that produced them.

\subsection{Required model-quality comparisons}

Evaluation uses a fixed tokenizer and held-out data for dense-versus-QAT
comparisons. Each of the seven Indic languages receives its own language
loss and task score, with a macro view that does not conceal low-resource
regressions. The scorecard also covers English, code, mathematics, native
script, transliteration where relevant, and code-mixed instructions.
Downstream tasks include understanding, instruction following, factual
answers, generation, and representative offline document workflows. Safety
and refusal behavior are assessed by language. Every task result records
dataset version, prompt, scoring method, decoding settings, checkpoint,
tokenizer, and contamination review.

Quantization is evaluated at three points: the dense parent, the immediate
ternary conversion before an optimizer update, and the final QAT checkpoint.
The same held-out weighting is applied to each. A matched dense continuation
provides a control for extra training and schedule changes. Report both
per-token likelihood and task outcomes; a lower weight-reconstruction error
or higher ternary zero fraction is not by itself a language-quality result.

\subsection{Systems and device evaluation}

Systems reporting uses valid tokens divided by elapsed time across all
ranks, plus slowest-rank step time, p50/p95 update latency, allocated and
reserved GPU memory, data waits, exposed communication, and scheduled
checkpoint and evaluation costs. Exact-resume tests compare model,
optimizer, RNG, and next-data position with an uninterrupted control.
The dense and QAT phases have separate speed and quality accounting.

For offline deployment, report exported weight-file bytes, scale metadata,
loaded peak RAM, KV-cache growth, prefill and decode latency, energy use,
and device compatibility on actual target hardware. No speculative drafter,
special ``thinking mode,'' Android-specific ternary kernel, or mobile speed
gain is part of the selected training model or reported as implemented here.

\section{Ternary storage and deployment limits}

With simple two-bit packing, the 1,349,222,400 eligible projection weights
need 337,305,600 bytes. A BF16 copy of the 150,994,944 tied-table weights
needs 301,989,888 bytes; the other 239,152 parameters add 478,304 BF16
bytes. Thus the arithmetic lower bound for that representation is
\begin{equation}
337{,}305{,}600+301{,}989{,}888+478{,}304
=639{,}773{,}792\ \text{bytes}
\end{equation}
\emph{before} ternary scales, file metadata, alignment, activations, and
runtime state. A fixed-rate base-3 code at $\log_2(3)$ bits per eligible
weight lowers the projection term but still does not establish device RAM.
Quantizing the tied table would require its own accuracy and storage
evaluation. A sub-700 MB loaded-RAM claim cannot be inferred from the
two-bit weight arithmetic and is not made here.

The training system and exported inference system have different storage
requirements. A packed inference artifact cannot be used to estimate
training-state memory, and a training throughput receipt cannot predict
on-device token generation. The release record must include the exact
artifact format, export parity against the trained QAT model, and device
measurements before describing the model as deployable on a given memory
class.

\section{Limitations and conclusion}

The selected geometry, dense optimizer, and computational implementation
are specified and supported by smaller-scale research and short single-GPU
measurements. They do not prove one-trillion-token feasibility within the
planning GPU-hour envelope, full-scale dense stability, preservation of
quality after ternary QAT, useful performance in every target language,
or offline device readiness. The current quantizer uses a coarse
128-by-128 scale tile; the best ternary grouping and QAT schedule at
1.5B have not been measured. Training data mixture and phase token
allocation are bound by versioned external run manifests rather than by
the model definition.

\model{} is a concrete dense-to-ternary design: a 25-layer,
1.500B-parameter transformer with Peri-LN and a learned cross-layer value
residual, trained first with FP8 arithmetic and then evaluated through
ternary continuation. The one-trillion-token program and seven-language
Indic scope define the intended experiment and reporting boundaries.
Final claims of language capability, compression, runtime memory, and
compute efficiency follow only from the completed checkpoints and device
measurements.

\begin{thebibliography}{10}
\bibitem{rmsnorm} B. Zhang and R. Sennrich. Root Mean Square Layer Normalization. \emph{NeurIPS}, 2019. \url{https://arxiv.org/abs/1910.07467}.
\bibitem{swiglu} N. Shazeer. GLU Variants Improve Transformer. 2020. \url{https://arxiv.org/abs/2002.05202}.
\bibitem{gqa} J. Ainslie et al. GQA: Training Generalized Multi-Query Transformer Models from Multi-Head Checkpoints. 2023. \url{https://arxiv.org/abs/2305.13245}.
\bibitem{bitnet} S. Ma et al. The Era of 1-bit LLMs: All Large Language Models are in 1.58 Bits. 2024. \url{https://arxiv.org/abs/2402.17764}.
\bibitem{fa3} J. Shah et al. FlashAttention-3: Fast and Accurate Attention with Asynchrony and Low-precision. 2024. \url{https://arxiv.org/abs/2407.08608}.
\bibitem{peri} J. Kim et al. Peri-LN: Revisiting Normalization Layer in the Transformer Architecture. 2025. \url{https://arxiv.org/abs/2502.02732}.
\bibitem{valueresidual} Z. Zhou et al. Value Residual Learning. 2024. \url{https://arxiv.org/abs/2410.17897}.
\bibitem{muon} J. Liu et al. Muon is Scalable for LLM Training. 2025. \url{https://arxiv.org/abs/2502.16982}.
\bibitem{deepgemm} DeepGEMM: Clean and Efficient FP8 GEMM Kernels with Fine-grained Scaling. \url{https://github.com/deepseek-ai/DeepGEMM}.
\end{thebibliography}

\end{document}
